\subsection{W-6 (BIG)}

Typické databázové stroje pro Big Data (Cassandra nebo HBase) - architektura, databázový model, distribuce dat.

\subsubsection*{Cassandra}

Všeobecné vlastnosti CassandraDB:

\begin{itemize}
  \item Data ukládána po sloupcích.
  \item Distribuovaná databáze. 
  \item Z Dynama si bere principy pro replikaci a distribuci dat:
  \begin{itemize}
    \item Dělení datovýsch setů pomocí konzistentního hashování
    \item Replikace dat pomocí verzovaných dat.
    \item Detekce a oprava chyb pomocí protokolu Gossip.
  \end{itemize}
  \item Vlastní dotazovací jazyk CQL (Cassandra Query Language) (neumí joiny a subqueries).
  \item Vysoká škálovatelnost - automaticky se vyvažuje i opravuje.
\end{itemize}

Výhody a nevýhody ukládání dat po sloupcích:
\begin{itemize}
  \item Výhody:
  \begin{itemize}
    \item Efektivní a rychlé agregační funkce.
    \item Komprese dat po sloupcích - ukládají se pouze existující data, ne null.
    \item Rychlý přístup k většímu množství dat.
  \end{itemize}
  \item Nevýhody:
  \begin{itemize}
    \item Nízký výkon pro úpravy dat - nutná úprava všech sloupců, mazání dat může být také pomalé, protože je potřeba opět smazat data v každém sloupci.
    \item Nižší výkon pro výběr omezeného množství malých dat.
    \item Méně flexibilní pro složitější dotazy (je to prostě tupá databáze).
  \end{itemize}
\end{itemize}

Infrastruktura CassandraDB:
\begin{itemize}
  \item Všechny uzly jsou si rovny (peer-to-peer).
  \item Cluster < Datacenter (machine) < Node.
\end{itemize}

CAP theorem:
Defaultně běží v AP módu, tedy Availability a Partition tolerance. Nesplňuje požadavky na Consistency - data se sice zapíšou, ale není zaručeno kdy. Cassandru lze nakonfigurovat tak, aby splňovala CP.

Optimalizace přístupu k datům:
\begin{itemize}
  \item Commit log - obdoba WAL, zapisuje se do něj každá změna a je používán pro obnovu.
  \item Memtable - datová struktura v paměti, do které se data zapisují po zápisu do commit logu.
  \item SSTable - diskový soubor, do kterého se flushují data z memtable jakmile začne přetékat. Při zápisu z memtable do SSTable se propsaná data odebírají z commit logu.
  \item Bloom filter - rychlý a nedeterministický vyhledávač, který určuje, zda je hodnota v SSTable.
\end{itemize}

Programovatelná konzistence:
Existují tři míry konzistence:
\begin{itemize}
  \item ONE - stačí přečíst jednu repliku. 
  \item QUORUM - stačí přečíst nadpoloviční většinu replik.
  \item ALL - musí se přečíst všechny repliky.
\end{itemize}

Architektura dělení dat:
\begin{itemize}
  \item Data se do uzlů dělí na základě Partition Key - součást primárního klíče, definuje na který uzel se data uloží.
  \item Ve vybraném uzlu poté určí fyzickou adresu dat Clustering Key.
  \item Data jsou v uzlech uložena do kruhu - uložená data se navzájem mohou překrývat (uzel může pokrýt více "stupňů v kruhu").
\end{itemize}
